\chapter{Phase 0 --- Setting Up Your Research Environment}\label{ch:setup}

\weeks{1--2} \quad Milestone M0

Phase~0 builds everything the rest of the project stands on: a repository with
a sensible structure, a working hardware toolchain, evaluation data, a trusted
metrics harness, and a measured \emph{float baseline}. Every later result is
compared against that baseline. It is unglamorous work, but a day saved here
costs a week later.

\begin{whybox}
RQ1 and RQ2 are questions about \emph{differences}: fixed point against float,
\qty{8}{\ms} against \qty{21}{\ms}. A difference is only meaningful if the
reference is measured carefully and reproducibly, with the same data, the same
metric code and the same alignment. That reference is the thesis pipeline, so
Phase~0 measures it properly for the first time on a standard benchmark.
\end{whybox}

\begin{center}\small
\begin{tabularx}{\textwidth}{@{}LL@{}}
\toprule
\textbf{Inputs} & \textbf{Outputs} \\
\midrule
The thesis repository; a computer (macOS, Linux or Windows); about 30\,GB of
free disk space; access to an x86-64 Linux or Windows machine for Vivado later
& \file{intent-se-hw} repository with CI; working toolchain; VoiceBank+DEMAND on
disk; \file{metrics/quality.py} with tests; baseline results table; lab notebook
\\
\bottomrule
\end{tabularx}
\end{center}

% =========================================================================
\section{Steps}\label{sec:setup:steps}

\begin{stepbox}[label=st:setup:notebook]{Start a lab notebook}
Create \file{docs/notebook/} in the new repository (next step) and write one
Markdown file per working day, named \file{YYYY-MM-DD.md}. Each entry records
the goal, what you did, the commands you ran, the commit hash, the results and
the next step. A template is in \cref{app:templates}.

Why so formal? In week 14 you will need to know exactly which configuration
produced the plot in week 5. Memory will not tell you, and your notebook will.

\textbf{Check:} today's entry exists and says what Phase~0 must achieve.
\textbf{Effort:} 30 minutes now, then 5 minutes per session for the rest of the
project.
\end{stepbox}

\begin{stepbox}[label=st:setup:repo]{Create the \texttt{intent-se-hw} repository}
Keep the thesis repository as an honest record of the thesis, and start a new
one for the hardware project. It imports the thesis package as its
floating-point reference.

\begin{lstlisting}[style=plain]
intent-se-hw/
  golden/        bit-accurate Python model (Phase 1)
  metrics/       PESQ, STOI, SI-SDR, delta-SNR
  experiments/   one folder per experiment: config.yaml, results.csv, notes.md
  nlp_hw/        classifier-head quantisation, encoder studies
  rtl/           SystemVerilog (Phase 2)
  tb/            cocotb testbenches (Phase 3)
  fpga/pynq_z2/  Vivado Tcl scripts, overlay, notebooks (Phase 4)
  asic/          LibreLane configuration, constraints, reports (Phase 5B)
  docs/          spec, notebook, figures
  .github/workflows/  CI
\end{lstlisting}

\begin{lstlisting}[style=shell]
mkdir intent-se-hw && cd intent-se-hw && git init
python3 -m venv .venv && source .venv/bin/activate
# the thesis code as an editable dependency (adjust the path)
pip install -e ../Intention_based_speech_enhancment_using_NLP
pip install numpy scipy soundfile pandas matplotlib pytest ruff pesq pystoi pyyaml
pip freeze > requirements-lock.txt
\end{lstlisting}

Add a \file{pyproject.toml} for your own packages (\file{golden},
\file{metrics}, \file{nlp_hw}) and a minimal \file{README.md} describing the
goal and layout. Commit.

\textbf{Check:} \code{python -c "from intent\_se.audio import SpeechEnhancer;
print(SpeechEnhancer().self\_test())"} prints a value around \num{4e-16}.
\textbf{Effort:} 1--2 hours.
\end{stepbox}

\begin{stepbox}[label=st:setup:tools]{Install the open-source hardware toolchain}
You need a simulator and linter (\tool{Verilator}~\cite{verilator}), a
Python testbench framework (\tool{cocotb}~\cite{cocotb}), a second simulator as
a cross-check (\tool{Icarus Verilog}), a waveform viewer (\tool{GTKWave} or
\tool{Surfer}), and a synthesis tool for quick area estimates
(\tool{Yosys}~\cite{yosys}).

The easiest route on Linux and macOS is the \textbf{YosysHQ OSS CAD Suite}. It
is a single download that bundles Yosys, Verilator, Icarus Verilog, GTKWave,
SymbiYosys and more. Unpack it and source its environment script. On macOS,
Homebrew (\code{brew install verilator icarus-verilog yosys}) also works.
Distribution packages of Verilator are often too old for current cocotb, so
check the minimum Verilator version in the cocotb documentation.

Install cocotb into your project's virtual environment so it uses the same
Python as your golden model: \code{pip install cocotb}.

Smoke-test everything with a ten-line counter:

\begin{lstlisting}[style=sv,caption={\texttt{sandbox/counter.sv}}]
module counter (input logic clk, input logic rst, output logic [7:0] count);
  always_ff @(posedge clk) begin
    if (rst) count <= '0;
    else     count <= count + 8'd1;
  end
endmodule
\end{lstlisting}

\begin{lstlisting}[style=py,caption={\texttt{sandbox/test\_counter.py}}]
import cocotb
from cocotb.clock import Clock
from cocotb.triggers import ClockCycles

@cocotb.test()
async def counts_up(dut):
    cocotb.start_soon(Clock(dut.clk, 10, "ns").start())
    dut.rst.value = 1
    await ClockCycles(dut.clk, 3)
    dut.rst.value = 0
    await ClockCycles(dut.clk, 2)
    start = int(dut.count.value)
    await ClockCycles(dut.clk, 5)
    assert int(dut.count.value) == start + 5
\end{lstlisting}

\begin{lstlisting}[style=shell]
verilator --version && yosys -V && iverilog -V | head -1
verilator --lint-only -Wall sandbox/counter.sv          # no warnings expected
yosys -p "read_verilog -sv sandbox/counter.sv; synth -top counter; stat"
\end{lstlisting}

Run the cocotb test with cocotb's Makefile flow or its Python runner. The
variable names changed between cocotb~1.x and~2.x (for example \code{TOPLEVEL}
became \code{COCOTB\_TOPLEVEL}), so copy the quick-start example for your
installed version from the cocotb documentation~\cite{ch12-13:cocotbdocs}.

\textbf{Check:} the lint is clean, Yosys reports a cell count, and the cocotb
test passes on Verilator \emph{and} on Icarus. \textbf{Effort:} 3--5 hours.
Tool installation is where newcomers lose time, so timebox it and ask for help
early.
\end{stepbox}

\begin{stepbox}[label=st:setup:vivado]{Secure a Vivado host now (you need it in week 12)}
AMD Vivado, needed for the PYNQ-Z2 in Phase~4, runs on x86-64 Linux and Windows
only~\cite{ch12-13:ug973}. It does not run on macOS, and Apple-silicon Macs
cannot run it natively. If your main computer is a Mac, arrange one of these
\emph{now}:
\begin{itemize}
  \item a lab or project PC at your university;
  \item a remote Linux machine you can reach by SSH or remote desktop;
  \item a cloud x86-64 Linux VM. Vivado needs several CPU cores, plenty of RAM
        and a large disk.
\end{itemize}
The free edition of Vivado supports the Zynq-7020 on the PYNQ-Z2. Check which
Vivado version your chosen PYNQ image release expects before downloading
(\cref{ch:fpga}); the installer is very large.

\textbf{Check:} you can log in to the machine and have enough disk space for
Vivado. \textbf{Effort:} 1--2 hours to arrange; the installation itself can
happen in the background later.
\end{stepbox}

\begin{stepbox}[label=st:setup:data]{Get the evaluation data}
Download \textbf{VoiceBank+DEMAND}~\cite{valentini2016voicebank,valentini2017dataset}
from the University of Edinburgh DataShare (DOI \code{10.7488/ds/2117}). It is the
most widely used benchmark for single-channel speech enhancement. It is recorded
at \qty{48}{\kHz}, the same rate as your thesis system, and it provides clean and
noisy versions of every utterance.

\begin{itemize}
  \item \textbf{Test set:} 824 utterances from two unseen speakers, at SNRs of
        2.5, 7.5, 12.5 and \qty{17.5}{\dB}. This is where you report results.
  \item \textbf{Training set:} 11,572 utterances from 28 speakers, at SNRs of 0,
        5, 10 and \qty{15}{\dB}. You do not train anything on the audio side, but
        use a fixed subset of it as a \emph{development} set for tuning (LUT sizes,
        word lengths), so the test set is touched only once.
        This is the same discipline as your thesis's classifier
        selection.
\end{itemize}
Also generate a small set of \textbf{synthetic signals} for hardware testing:
silence, a full-scale sine, a quiet sine in white noise, an impulse, a chirp, DC,
and clipped noise. Real speech hides bugs that these signals expose
immediately (\cref{ch:verification}).

\textbf{Check:} a script counts 824 clean and 824 noisy test files, all at
\qty{48}{\kHz}, and the synthetic signals are generated deterministically from a
fixed seed. \textbf{Effort:} 2 hours plus download time.
\end{stepbox}

\begin{stepbox}[label=st:setup:metrics]{Build and test the metrics harness}
Write \file{metrics/quality.py} with four metrics. Each one takes the clean
reference and the processed signal, already aligned
(\cref{st:setup:baseline}).

\begin{itemize}
  \item \textbf{PESQ}, wideband, ITU-T P.862.2~\cite{itu2001pesq,itu2007pesqwb}.
  The common Python package \code{pesq} accepts only 8 or \qty{16}{\kHz}
  signals, so resample from 48 to \qty{16}{\kHz} first.
  \item \textbf{STOI}~\cite{taal2011stoi}, via \code{pystoi}.
  \item \textbf{SI-SDR}~\cite{leroux2019sisdr}, in a few lines of numpy. It has
  no dependencies, so it is your fallback if a package breaks.
  \item \textbf{$\Delta$SNR}: the output SNR minus the input SNR, both against
  the clean reference.
\end{itemize}

\begin{lstlisting}[style=py,caption={The core of \texttt{metrics/quality.py}.}]
import numpy as np
from math import gcd
from scipy.signal import resample_poly
from pesq import pesq
from pystoi import stoi

def to_16k(x, fs):
    g = gcd(16000, fs)
    return resample_poly(x, 16000 // g, fs // g)

def si_sdr(ref, est):
    ref = ref - ref.mean(); est = est - est.mean()
    target = (np.dot(est, ref) / np.dot(ref, ref)) * ref
    noise = est - target
    return 10 * np.log10(np.sum(target**2) / np.sum(noise**2))

def quality(ref, est, fs=48000):
    return {
        "pesq_wb": pesq(16000, to_16k(ref, fs), to_16k(est, fs), "wb"),
        "stoi": stoi(ref, est, fs, extended=False),
        "si_sdr": si_sdr(ref, est),
    }
\end{lstlisting}

Write unit tests (\file{metrics/test_quality.py}) that pin down each metric's
behaviour:
\begin{itemize}
  \item identical signals give PESQ near its maximum, STOI near 1, and a very
        high SI-SDR;
  \item a clean signal plus white noise at a known SNR gives an SI-SDR within
        about \qty{0.1}{\dB} of that SNR;
  \item scaling the estimate leaves SI-SDR unchanged, since that is the
        \enquote{scale-invariant} part.
\end{itemize}

\textbf{Check:} \code{pytest metrics} passes. \textbf{Effort:} 3--4 hours.
\end{stepbox}

\begin{stepbox}[label=st:setup:baseline]{Measure the float baseline}
Run the thesis \code{SpeechEnhancer} over the whole test set and record the
metrics per utterance. Do this for the unprocessed noisy input and for
$\beta\in\{0.5, 1.0, 1.5, 2.0\}$ with the other parameters at their defaults.

Two details matter:
\begin{enumerate}
  \item \textbf{Alignment.} Offline, \code{process\_signal} returns output
  delayed by exactly 512 samples (one hop) relative to its input. This was
  measured with an impulse. Shift the output by 512 samples before computing
  metrics, otherwise PESQ and SI-SDR punish a delay that is not a quality
  problem. Verify the shift yourself by cross-correlation on one file.
  \item \textbf{Pre-roll.} The first 5 frames (about \qty{53}{\ms}) pass through
  unprocessed while IMCRA is seeded. Keep this behaviour, since the hardware
  will have it too, but note it in your results.
\end{enumerate}

Save everything under \file{experiments/000-baseline/}:
\begin{itemize}
  \item \file{config.yaml}: parameters, commit hash, package versions;
  \item \file{results.csv}: one row per utterance and setting;
  \item \file{summary.md}: means and 95\% confidence intervals per setting and
        per input SNR (\cref{ch:evaluation}).
\end{itemize}

\textbf{Check:} enhancement improves SI-SDR and $\Delta$SNR over the noisy
input. If PESQ or STOI get worse at some $\beta$ values, do not panic: it is
a known property of Wiener-type filters, and a finding worth writing down. The
whole run is reproducible from one command. \textbf{Effort:} 4--6 hours
(the run itself takes minutes).
\end{stepbox}

\begin{stepbox}[label=st:setup:ci]{Set up continuous integration}
Add a GitHub Actions workflow (\file{.github/workflows/ci.yml}) that installs
the Python environment and runs \code{ruff} and \code{pytest} on every push.
In Phase~3 you will add Verilator and the cocotb regression to the same
workflow.

\textbf{Check:} the CI badge is green on the main branch.
\textbf{Effort:} 1--2 hours.
\end{stepbox}

% =========================================================================
\section{Pitfalls}\label{sec:setup:pitfalls}

\begin{pitfallbox}[Pitfalls in Phase 0]
\begin{itemize}
  \item \textbf{Unaligned signals.} Forgetting the 512-sample delay makes every
        metric look worse than it is, and the error changes with the frame
        configuration in Phase~1, which corrupts your comparisons.
  \item \textbf{Tuning on the test set.} Choose every parameter on the
        development subset. Touch the test set only to report results.
  \item \textbf{Unpinned versions.} PESQ implementations and resamplers differ
        slightly. Pin package versions and record them with each result.
  \item \textbf{Leaving Vivado until week 12.} A Mac cannot run it. Solve this
        now.
  \item \textbf{Polishing the setup forever.} Phase~0 has two weeks. A working,
        tested harness is enough; elegance can come later.
\end{itemize}
\end{pitfallbox}

\begin{donebox}[Milestone M0 is reached when]
\begin{checklist}
  \item \file{intent-se-hw} exists with the agreed layout, a lock file and green CI.
  \item Verilator, Icarus, Yosys, cocotb and a waveform viewer pass their smoke tests.
  \item A Vivado host is arranged.
  \item VoiceBank+DEMAND and the synthetic signals are on disk.
  \item \file{metrics/quality.py} is implemented and unit-tested.
  \item The float baseline is measured, aligned, saved and summarised.
  \item The lab notebook has an entry for every working day so far.
\end{checklist}
\end{donebox}
